\scanpage{121}
une sorte de synthèse géométrico-arithmétique des nombreuses théories cohomologiques dont on dispose\deleted{\unsure{rait}} à présent pour les variétés algébri\deleted{\unsure{qu}}ques,\deleted{de toutes dimensions,}et des représentations linéaires de groupes de \deleted{g}Galois arithmétiques auxquelles ces théories donnent naissance\deleted{\unsure{s}}.

5) \uline{Solution du problème d'introduction d'une} $\handwritten{\lambda}$-\uline{structure} (correspondant à la notion de puissance extérieure) \uline{dans l'anneau} $K^\bullet(X)$(dit ``de Grothendieck'')d'un espace topologique localement annelé (par exemple une variété analytique, ou\deleted{\unsure{nd|md}} \unsure{différentiable}, ou un schéma)\handwritten{,} construit à l'aide de la catégorie des complexes ``parfaits'' sur $X$. Ce problème était signalé comme le plus important problème \inserted{\handwritten{resté}} en suspens dans le Séminaire de Géomé\deleted{y}\handwritten{t}rie Algébrique de l'IHES 1966/67, sur le théorème de Riemann-Roch. La nécessité de travailler avec le $K^\bullet(X)$ défini en termes de \uline{complexes} parfaits, plutôt que de \uline{fibrés} vectoriels seulement, est \unsure{imposé|imposée} par des \deleted{nécess} impératifs techniques irrécusables, qui finiront certainement par s'imposer aussi en dehors de la \deleted{\unsure{G}}\handwritten{g}éométrie algébrique. La solution de Deligne, comme prévu, passe par l'intermédiaire d'une construction de \deleted{\unsure{puis}} foncteurs puissances extérieures dans la catégorie dérivée $D(X)$ (catégorie de\deleted{s} complexes de modules \inserted{\handwritten{sur $X$}}, - en fait, il faut se borner ici aux complexes de \uline{chaines}), en utilisant la théorie de Dold\inserted{\handwritten{-Puppe}}. La difficulté sérieuse consistait à prouver que ces opérations dans $K^\bullet(X)$ satisfont aux propriétés algébriques habituelles, permettant d'exprimer par exemple $\handwritten{\lambda}^i(xy)$ comme polynôme universel à coéfficients entiers en les $\handwritten{\lambda}^j(x),\handwritten{\lambda}^k(y)$. Deligne procède par une voie très ingénieuse, en déterminant d'abord la structure \unsure{da|du} $\handwritten{\lambda}$-anneau de toutes les ``opérations tensorielles universelles'' qu'on peut effectuer sur des fibrés vectoriels. \deleted{\unsure{E}}Il n'a écrit sa théorie que lorsqu'on travaille au dessus d'un corps de base \deleted{\unsure{finis|fini|choisi}}(de caractéristique arbitraire). \deleted{Le cas de la caractéristique}
% archive sheet number: 120
